\section{Proofs and Details for Composition}
\label{app:composition}

This appendix proves the results of Section~\ref{sec:composition}, defines
composition with independent inputs and its curried evaluation, shows that the
converse of Corollary~\ref{cor:monotone-commitment} fails, and works out the
replicated store with bare reads.

\subsection{Sequential composition}

\begin{proof}[Proof of Lemma~\ref{lem:composition}]
We verify the common-\indesc condition
(Proposition~\ref{prop:common-expansion-form}).  Let rows $(i_r,o_r)$ of
$S_2\circ S_1$, finitely many, have common \indesc $j$, and choose $m_r$ with
$(i_r,m_r)\in S_1$ and $(m_r,o_r)\in S_2$.  The condition for $\Spec_1$ gives
$m\in S_1(j)$ \outrelpart every $m_r$.  Then $m$ is a common \indesc of the
inputs of the rows $(m_r,o_r)$, and the condition for $\Spec_2$ gives
$u\in S_2(m)$ \outrelpart every $o_r$.  So $(j,u)\in S_2\circ S_1$ is the
required common \outnoun.
\end{proof}

Tables of functions are closed under $\circ$, so Lemma~\ref{lem:composition}
specializes to functions.

\begin{corollary}[Functional composition]
\label{cor:functional-composition}
If $f:(I,\InputOrd)\to(M,\preceq_M)$ and
$g:(M,\preceq_M)\to(O,\OutputOrd)$ are monotone, so is $g\circ f$.
\end{corollary}

\begin{proof}
The tables of $f$ and $g$ satisfy Lifted Monotonicity and compose to the
table of $g\circ f$; apply Lemma~\ref{lem:composition}.
\end{proof}

\subsection{Composition with independent inputs}
\label{sec:parameterized-composition}

Sequential composition cannot express a downstream component with input of
its own, such as a replica that serves bare reads
(Section~\ref{sec:composition}), so we give each component a separate
input.  Let
\begin{align*}
  \Spec_F&=\bigl(I\times A,\preceq_I\times\preceq_A,B,\preceq_B,F\bigr),\\
  \Spec_G&=\bigl(J\times B,\preceq_J\times\preceq_B,C,\preceq_C,G\bigr)
\end{align*}
be ordered relational specifications sharing the ordered domain
$(B,\preceq_B)$, which carries values from $F$ to $G$.  Here $i\in I$ is input
that only $F$ receives, and $j\in J$ is input that only $G$ receives.  We
split $F$'s input as $(i,a)$ so that $F$ restricted to a fixed $i$ is a table
from $A$ to $B$.  The curried evaluation below uses this split to evaluate
the composition one pair $(i,j)$ at a time.  The \emph{independently
parameterized composition} of $\Spec_F$ and $\Spec_G$ is
\begin{align}
  G\star F
  &=\pi_{I,J,A,C}\bigl(F\Join_{B}G\bigr) \notag\\
  &=\{((i,j,a),c):\exists b.\ ((i,a),b)\in F
                         \land ((j,b),c)\in G\},
  \label{eq:parameterized-composition}
\end{align}
with the product order on $I\times J\times A$.  Since $F$ is total on
$I\times A$ and $G$ is total on $J\times B$, every composed input has a
nonempty fiber.

\begin{lemma}[Independently parameterized composition]
\label{lem:parameterized-composition}
If $\Spec_F$ and $\Spec_G$ satisfy Lifted Monotonicity, then the specification
with table $G\star F$ satisfies Lifted Monotonicity.
\end{lemma}

\begin{proof}
Let rows $((i_r,j_r,a_r),c_r)$, finitely many, have common \indesc $(i,j,a)$,
with witnesses $b_r$ from Equation~\ref{eq:parameterized-composition}.  The
common-\indesc condition for $\Spec_F$ at $(i,a)$ gives $b\in F(i,a)$
\outrelpart every $b_r$.  Then $(j,b)$ is a common \indesc of the
$(j_r,b_r)$, and the condition for $\Spec_G$ gives $c\in G(j,b)$ \outrelpart
every $c_r$.  So $((i,j,a),c)\in G\star F$ is the required common \outnoun.
\end{proof}

\paragraph{Curried evaluation.}
For fixed parameters let
\begin{equation}
  F_i=\pi_{A,B}(F[I=i]),
  \qquad
  G_j=\pi_{B,C}(G[J=j]),
  \label{eq:curried-fibers}
\end{equation}
with the selection of Section~\ref{sec:functional-specifications}.

\begin{proposition}[Curried equijoin plan]
\label{prop:curried-query-plan}
The table in Equation~\ref{eq:parameterized-composition} equals
\begin{equation}
  \bigcup_{(i,j)\in I\times J}
  \left\{((i,j,a),c):
    (a,c)\in\pi_{A,C}\bigl(F_i\Join_B G_j\bigr)\right\}.
  \label{eq:curried-query-plan}
\end{equation}
\end{proposition}

\begin{proof}
By Equation~\ref{eq:curried-fibers}, $((i,a),b)\in F$ and $((j,b),c)\in G$
iff $(a,b)\in F_i$ and $(b,c)\in G_j$.  Quantifying over $b$ gives the same
condition on $((i,j,a),c)$ in both expressions.
\end{proof}

For monotone $f:I\times A\to B$ and $g:J\times B\to C$, the instance of
Lemma~\ref{lem:parameterized-composition} for their tables says that
\begin{equation}
  h(i,j,a)=g\bigl(j,f(i,a)\bigr)
  \label{eq:curried-composition}
\end{equation}
is monotone; in curried form, $h^\sharp(i,j)=g^\sharp(j)\circ f^\sharp(i)$,
which is the plan of Proposition~\ref{prop:curried-query-plan} on tables of functions.

\subsection{Proper Coordination}

Definition~\ref{def:proper-coordination} and
Corollaries~\ref{thm:proper-coordination} and~\ref{cor:monotone-commitment}
apply to independently parameterized composition, with $F$ in place of $S_1$,
$G$ in place of $S_2$, and $G\star K$ as the resolved table.

\begin{proof}[Proof of Corollary~\ref{thm:proper-coordination}]
If the resolved specification satisfies Lifted Monotonicity, Complete CALM
gives consistency and Proposition~\ref{prop:exposure-boundary} gives the
exposability clause.  Conversely, properness includes consistency, and Complete
CALM gives Lifted Monotonicity.
\end{proof}

\begin{proof}[Proof of Corollary~\ref{cor:monotone-commitment}]
The table of a monotone function satisfies Lifted Monotonicity.  Apply
Lemma~\ref{lem:composition} or~\ref{lem:parameterized-composition}, then
Corollary~\ref{thm:proper-coordination}.
\end{proof}

The converse of Corollary~\ref{cor:monotone-commitment} fails: let
$I=\{*\}$, $M=\{x,y\}$ discrete, $S_1=\{(*,x),(*,y)\}$, $K=\{(*,x)\}$, and
$S_2=\{(x,x),(y,x),(y,y)\}$.  The fiber $S_2(y)=\{x,y\}$ is not directed, so
$S_2$ fails Lifted Monotonicity and the corollary does not apply, but the
resolved table $\{(*,x)\}$ is monotone, so the strategy is proper.

\begin{example}[Bare reads]
\label{ex:bare-reads}
We add bare reads to Example~\ref{ex:committed-log-composition}.  Take $I$
and $B$ the proposal streams and logs of that example, $A$ a one-point set,
and $J$ the histories of bare reads received by a replica, ordered by prefix.
Take $C$ the pairs of a response history for logged commands, ordered by
extension, and a set of bare-read responses, ordered by inclusion.  Let $F$ be
the log specification $S_1$ of that example, with the one-point $A$ added to
its input, and $K$ the arrival-order commitment.  Let $G$ relate a bare-read history $j$ and a log $b$ to the
responses obtained by executing $b$, together with at least one response to
each bare read in $j$.  Responses are facts, so a set of them may answer one
read more than once.  As in Example~\ref{ex:committed-log-composition}, $F$
and $G\star F$ are non-monotone, and $K$ is the table of a monotone function.

Suppose first that each bare-read response records the log position it read
from, as in ``$x=5$ as of entry 17.''  A longer log never contradicts such a
response, so $G$ is monotone and bare reads need no coordination.  By
Corollary~\ref{cor:monotone-commitment} the log commitments are proper, and
$G\star K$ is consistent.  Suppose instead that the specification demands the
latest value, so that a read of $x$ must return the value of $x$ after all of
$b$.  Then a bare read can be contradicted by an entry it has not seen, $G$ is
not monotone, and neither is $G\star K$.  Bare reads then need coordination,
typically by going through the log.
\end{example}
